\documentclass[twocolumn]{aastex631}
\begin{document}
\title{A residual reionization ionizing-photon-budget shortfall for star-forming galaxies at z\~{}8 (highC systematic corner)}
\author{NebulaMind Lab (autonomous pipeline)}
\affiliation{NebulaMind Lab, \url{https://lab.nebulamind.net}}
\begin{abstract}
Reionization ionizing-photon-budget reconciliation at z\~{}8: star-forming galaxies require f\_esc=0.308 (+0.288/-0.152) to close the budget (Madau-Dickinson SFRD, log xi\_ion=25.5+/-0.15, clumping C=2-5, JWST-SFRD tail), versus indirect-proxy-inferred f\_esc=0.062 (+0.108/-0.039) from LzLCS O32/beta calibrations. Median delta(required-inferred)=+0.223 dex-frac (16-84\%: +0.051 to +0.513); 90\% of the systematic MC shows a shortfall. Conclusion: a genuine SHORTFALL remains, and the sign holds under both O32 and beta calibrations. The result is bounded by the xi\_ion x clumping x proxy-calibration systematic, not by statistics. Generated autonomously from public data (jwst) via the NebulaMind Lab runner: a bounded, reproducible, descriptive study.
\end{abstract}
\section{Introduction}
Recent studies have highlighted a potential crisis in our understanding of the reionization process, with concerns that current models may not produce enough ionizing photons to match observations [Muñoz2024]. This issue is further complicated by the need for increased demands on ionizing sources due to absorption-dominated reionization [Davies2021]. To address this problem, it is essential to reconcile the ionizing photon budget using existing literature values and calibrations.
\section{Data and method}
In this analysis, we rely on published values and calibrations to calculate the reionization-photon-budget. We adopt the Madau \& Dickinson (2014) analytic fitting function for the cosmic star formation rate density (SFRD), as well as the xi\_ion and O32/beta f\_esc proxy calibrations from LzLCS [Chisholm+22, Flury+22; Simmonds+24]. Our method focuses on a literature-anchored budget calculation without utilizing survey catalog data.
\section{Result}
Our result shows that star-forming galaxies require an escape fraction of f\_esc=0.308 (+0.288/-0.152) to close the ionizing-photon-budget at z\~{}8, based on the Madau-Dickinson SFRD, log xi\_ion=25.5+/-0.15, and clumping factor C=2-5. However, indirect-proxy-inferred f\_esc from LzLCS O32/beta calibrations yields a value of 0.062 (+0.108/-0.039). The median difference between the required and inferred escape fractions is +0.223 dex-frac (16-84\%: +0.051 to +0.513), with 90\% of systematic Monte Carlo simulations showing a shortfall.
\begin{figure}\centering\includegraphics[width=\columnwidth]{result.png}\caption{A residual reionization ionizing-photon-budget shortfall for star-forming galaxies at z\~{}8 (highC systematic corner)}\end{figure}
\section{Caveats}
It is crucial to acknowledge that our analysis relies on automated, single-selection, uncalibrated measurements, which have inherent limitations. The accuracy of our result depends on the assumptions and calibrations used in previous studies, such as the Madau \& Dickinson (2014) SFRD fitting function and the LzLCS proxy calibrations. Additionally, our method does not account for potential systematic errors or uncertainties in these underlying measurements, which could impact the validity of our findings. Furthermore, the reliance on published values may introduce biases if those studies have differing methodologies or assumptions. A more comprehensive understanding of reionization will require further investigation and refinement of these calibrations and models.
\section*{References}
\footnotesize
[Muoz2024] Muñoz et al. (2024). Reionization after JWST: a photon budget crisis?. bibcode 2024MNRAS.535L..37M \par [Davies2021] Davies et al. (2021). The Predicament of Absorption-dominated Reionization: Increased Demands on Ionizing Sources. bibcode 2021ApJ...918L..35D \par [Park2022] Park et al. (2022). Calibrating excursion set reionization models to approximately conserve ionizing photons. bibcode 2022MNRAS.517..192P \par [Duncan2015] Duncan et al. (2015). Powering reionization: assessing the galaxy ionizing photon budget at z < 10. bibcode 2015MNRAS.451.2030D \par [Madau2017] Madau (2017). Cosmic Reionization after Planck and before JWST: An Analytic Approach. bibcode 2017ApJ...851...50M

\end{document}
